% GENERATED FILE. Edit canonical lessons, then run npm run generate:books.
% canonical-source-hash: fnv1a64:cbca1792333ecfbe
% canonical-lessons: JA-C429-gakkari-suru, JA-C429-ochikomu, JA-C429-tameiki, JA-C429-tsurai, JA-C429-nagusameru

\chapter{Feeling Down}
\label{ch:ja-feeling-down}
\hlchaptermodality{429}

\begin{chapteropening}
\textbf{By the end of this chapter:} \emph{I can say I am disappointed and comfort a friend.}

Complete the last lesson of chapter 429: I can say I am disappointed and comfort a friend.
\end{chapteropening}

\section[\texorpdfstring{\ja{がっかりする} (gakkari suru)}{がっかりする (gakkari suru)}]{\ja{がっかりする} (gakkari suru) — to be disappointed}

\label{lesson:JA-C429-gakkari-suru}



\begin{quote}
Before the new one: say the Japanese for to frolic, then the Japanese for to be full of energy.
\end{quote}

\subsection*{You'll want to know: \ja{がっかりする}}
\textbf{\ja{がっかりする}} — \emph{gakkari suru} — \textquotedblleft{}to be disappointed\textquotedblright{}.

\begin{quote}
\textbf{\ja{しあいに} \ja{まけて、がっかりしました。}} — \emph{shiai ni makete, gakkari shimashita} — \textquotedblleft{}We lost the match, and I was disappointed.\textquotedblright{}
\end{quote}

Polite \textbf{\ja{がっかりします}}, \emph{gakkarishimasu}; a \ja{する} verb. Past (polite): \textbf{\ja{せんしゅが} \ja{がっかりしました。}}, \emph{senshu ga gakkari shimashita}.

\begin{cousinweb}
\textbf{\ja{がっかりする}} (\emph{gakkari suru}) is a native Japanese mimetic word, \emph{gakkari}, inherited rather than borrowed from Chinese, made a verb with \emph{suru}, \textquotedblleft{}to do\textquotedblright{}; its older origin is not settled (JMdict).
\end{cousinweb}

\begin{grammarlens}[title={What you've built}]
One more word for describing what you lived and felt.
\end{grammarlens}

\subsection*{Your turn}
\begin{itemize}
\raggedright
  \item \emph{Say it:} \emph{gakkari suru}
  \item \emph{Say it:} \emph{gakkari suru}, once more
  \item \emph{Say it:} \emph{shiai ni makete, gakkari shimashita} — We lost the match, and I was disappointed.
  \item \emph{Recall:} say the Japanese for to frolic, then the Japanese for to be full of energy, then say \emph{gakkari suru} again
\end{itemize}

\begin{tcolorbox}[breakable,colback=teal!4,colframe=teal!35!black,title={Before you move on}]
What does \ja{がっかりする} mean? (\textquotedblleft{}To be disappointed\textquotedblright{}.) Say it once more.
\end{tcolorbox}
\section[\texorpdfstring{\ja{おちこむ} (ochikomu)}{おちこむ (ochikomu)}]{\ja{おちこむ} (ochikomu) — to feel down; to be depressed}

\label{lesson:JA-C429-ochikomu}



\begin{quote}
Before the new one: say the Japanese for to be full of energy, then the Japanese for to be disappointed.
\end{quote}

\subsection*{You'll want to know: \ja{おちこむ}}
\textbf{\ja{おちこむ}} — \emph{ochikomu} — \textquotedblleft{}to feel down; to be depressed\textquotedblright{}.

\begin{quote}
\textbf{\ja{しけんの} \ja{あと、あには} \ja{おちこんでいました。}} — \emph{shiken no ato, ani wa ochikonde imashita} — \textquotedblleft{}After the exam my elder brother was feeling down.\textquotedblright{}
\end{quote}

Polite \textbf{\ja{おちこみます}}, \emph{ochikomimasu}; a godan verb. Past (polite): \textbf{\ja{あにが} \ja{おちこみました。}}, \emph{ani ga ochikomimashita}.

\begin{cousinweb}
\textbf{\ja{おちこむ}} (\emph{ochikomu}) is a native Japanese compound (wago), inherited rather than borrowed from Chinese: \emph{ochi}, \textquotedblleft{}falling (ochiru)\textquotedblright{} and \emph{komu}, \textquotedblleft{}going in\textquotedblright{} (JMdict, KANJIDIC2).
\end{cousinweb}

\begin{grammarlens}[title={What you've built}]
One more word for describing what you lived and felt.
\end{grammarlens}

\subsection*{Your turn}
\begin{itemize}
\raggedright
  \item \emph{Say it:} \emph{ochikomu}
  \item \emph{Say it:} \emph{ochikomu}, once more
  \item \emph{Say it:} \emph{shiken no ato, ani wa ochikonde imashita} — After the exam my elder brother was feeling down.
  \item \emph{Recall:} say the Japanese for to be full of energy, then the Japanese for to be disappointed, then say \emph{ochikomu} again
\end{itemize}

\begin{tcolorbox}[breakable,colback=teal!4,colframe=teal!35!black,title={Before you move on}]
What does \ja{おちこむ} mean? (\textquotedblleft{}To feel down; to be depressed\textquotedblright{}.) Say it once more.
\end{tcolorbox}
\section[\texorpdfstring{\ja{ためいき} (tameiki)}{ためいき (tameiki)}]{\ja{ためいき} (tameiki) — a sigh}

\label{lesson:JA-C429-tameiki}



\begin{quote}
Before the new one: say the Japanese for to be disappointed, then the Japanese for to feel down.
\end{quote}

\subsection*{You'll want to know: \ja{ためいき}}
\textbf{\ja{ためいき}} — \emph{tameiki} — \textquotedblleft{}a sigh\textquotedblright{}.

\begin{quote}
\textbf{\ja{ちちは} \ja{ながい} \ja{ためいきを} \ja{つきました。}} — \emph{chichi wa nagai tameiki o tsukimashita} — \textquotedblleft{}My father let out a long sigh.\textquotedblright{}
\end{quote}

\emph{tameiki o tsuku} is \textquotedblleft{}to let out a sigh\textquotedblright{}: this \emph{tsuku} goes with \emph{tameiki} and is not \emph{tsuku}, \textquotedblleft{}to arrive\textquotedblright{}.

\begin{cousinweb}
\textbf{\ja{ためいき}} (\emph{tameiki}) is a native Japanese word (wago), inherited rather than borrowed from Chinese; its kanji spelling was attached to it later (JMdict).
\end{cousinweb}

\begin{grammarlens}[title={What you've built}]
One more word for describing what you lived and felt.
\end{grammarlens}

\subsection*{Your turn}
\begin{itemize}
\raggedright
  \item \emph{Say it:} \emph{tameiki}
  \item \emph{Say it:} \emph{tameiki}, once more
  \item \emph{Say it:} \emph{chichi wa nagai tameiki o tsukimashita} — My father let out a long sigh.
  \item \emph{Recall:} say the Japanese for to be disappointed, then the Japanese for to feel down, then say \emph{tameiki} again
\end{itemize}

\begin{tcolorbox}[breakable,colback=teal!4,colframe=teal!35!black,title={Before you move on}]
What does \ja{ためいき} mean? (\textquotedblleft{}A sigh\textquotedblright{}.) Say it once more.
\end{tcolorbox}
\section[\texorpdfstring{\ja{つらい} (tsurai)}{つらい (tsurai)}]{\ja{つらい} (tsurai) — hard to bear; painful}

\label{lesson:JA-C429-tsurai}



\begin{quote}
Before the new one: say the Japanese for to feel down, then the Japanese for a sigh.
\end{quote}

\subsection*{You'll want to know: \ja{つらい}}
\textbf{\ja{つらい}} — \emph{tsurai} — \textquotedblleft{}hard to bear; painful\textquotedblright{}.

\begin{quote}
\textbf{\ja{にゅういんは} \ja{つらかったです。}} — \emph{nyūin wa tsurakatta desu} — \textquotedblleft{}Being in hospital was hard to bear.\textquotedblright{}
\end{quote}

An i-adjective.

\begin{cousinweb}
\textbf{\ja{つらい}} (\emph{tsurai}) is a native Japanese word (wago), inherited rather than borrowed from Chinese; its kanji spelling was attached to it later (JMdict).
\end{cousinweb}

\begin{grammarlens}[title={What you've built}]
One more word for describing what you lived and felt.
\end{grammarlens}

\subsection*{Your turn}
\begin{itemize}
\raggedright
  \item \emph{Say it:} \emph{tsurai}
  \item \emph{Say it:} \emph{tsurai}, once more
  \item \emph{Say it:} \emph{nyūin wa tsurakatta desu} — Being in hospital was hard to bear.
  \item \emph{Recall:} say the Japanese for to feel down, then the Japanese for a sigh, then say \emph{tsurai} again
\end{itemize}

\begin{tcolorbox}[breakable,colback=teal!4,colframe=teal!35!black,title={Before you move on}]
What does \ja{つらい} mean? (\textquotedblleft{}Hard to bear; painful\textquotedblright{}.) Say it once more.
\end{tcolorbox}
\section[\texorpdfstring{\ja{なぐさめる} (nagusameru)}{なぐさめる (nagusameru)}]{\ja{なぐさめる} (nagusameru) — to comfort; to console}

\label{lesson:JA-C429-nagusameru}



\begin{quote}
Before the new one: say the Japanese for a sigh, then the Japanese for hard to bear.
\end{quote}

\subsection*{You'll want to know: \ja{なぐさめる}}
\textbf{\ja{なぐさめる}} — \emph{nagusameru} — \textquotedblleft{}to comfort; to console\textquotedblright{}.

\begin{quote}
\textbf{\ja{ともだちが} \ja{ないたので、なぐさめました。}} — \emph{tomodachi ga naita node, nagusamemashita} — \textquotedblleft{}My friend cried, so I comforted him.\textquotedblright{}
\end{quote}

Polite \textbf{\ja{なぐさめます}}, \emph{nagusamemasu}; an ichidan verb. Past (polite): \textbf{\ja{せんせいが} \ja{がくせいを} \ja{なぐさめました。}}, \emph{sensei ga gakusei o nagusamemashita}.

\begin{cousinweb}
\textbf{\ja{なぐさめる}} (\emph{nagusameru}) is a native Japanese word (wago), inherited rather than borrowed from Chinese; its kanji spelling was attached to it later (JMdict).
\end{cousinweb}

\begin{grammarlens}[title={What you've built}]
One more word for describing what you lived and felt.
\end{grammarlens}

\subsection*{Your turn}
\begin{itemize}
\raggedright
  \item \emph{Say it:} \emph{nagusameru}
  \item \emph{Say it:} \emph{nagusameru}, once more
  \item \emph{Say it:} \emph{tomodachi ga naita node, nagusamemashita} — My friend cried, so I comforted him.
  \item \emph{Recall:} say the Japanese for a sigh, then the Japanese for hard to bear, then say \emph{nagusameru} again
\end{itemize}

\begin{tcolorbox}[breakable,colback=teal!4,colframe=teal!35!black,title={Before you move on}]
What does \ja{なぐさめる} mean? (\textquotedblleft{}To comfort; to console\textquotedblright{}.) Say it once more.
\end{tcolorbox}
